\chapter{データの時点と先読み}\label{ch:pit}

金融の実証で最も多い誤りは、「その時点では知りえなかった情報」を使ってしまうことである。この章では、この研究で注意した先読みの経路と、その防ぎ方をまとめる。第 VI 部のデータ基盤は、この章の決まりを守るように作った。

\section{時点つきのデータ}

\begin{defbox}[定義：ポイント・イン・タイム]
各データに「いつから知られていたか」の時刻を付け、判断日 $d$ には、時刻が $d$ 以前のデータだけを使う。値が後から訂正されたときは、訂正の前の値も残す。
\end{defbox}

\section{先読みの経路と防ぎ方}

\begin{longtable}{P{3.0cm}P{5.4cm}P{6.0cm}}
\toprule
\textbf{データ} & \textbf{先読みの経路} & \textbf{この研究の防ぎ方} \\
\midrule
\endhead
\bottomrule
\endfoot
指数の構成銘柄 & 今の構成銘柄で過去を見る（生存者バイアス、後から入った銘柄の先取り） & Wikipedia の入れ替えの履歴から、各判断日の構成銘柄を作り直す \\
社名・ティッカー & 今の社名を LLM に渡すと、その後の社名の変更（合併など）が伝わる & 総合報告書 1 の時点では対応できず、限界として記録した \\
株価 & 上場を廃止した銘柄の株価がない。取得日の基準で調整された値を使う & 欠ける銘柄を記録する。上場をやめた銘柄は、その後のリターンを 0\% とする（実験 7〜9） \\
財務諸表 & 後の報告書で訂正・組み替えた値を、当時の値として使う & 同じ四半期に複数の値があれば、\textbf{最初に提出された値}を使う（実験 C で見つけた APA の不具合） \\
提出日 & 期末の日付で判断する（実際の提出は 1〜3 か月後） & 提出日（受理日）が判断日以前の決算だけを使う \\
決算発表文 & 受理日の終値に反映されたかを誤る & 受理時刻（UTC → 東部時間）と時間帯から day 0 を決める \\
アナリスト予想 & 発表の後に改訂された予想を使う & Alpha Vantage の「発表前のコンセンサス」を使う。集計の時点は公開されていない（限界） \\
月次レポート & 表紙の月で判断する & PDF の作成日を公表日とする。「$N$ 月号」は前月の下旬に作られていた \\
LLM の知識 & 学習の時点より前の出来事を、その後の結果ごと覚えている & gpt-4o-mini は知識の締め切り（2023-10）より後の期間だけ。ChronoGPT は前の年末の版 \\
学習したモデル & 評価する期間の反応で学習する & 四半期ごとのローリング、反応の窓が閉じた発表だけ \\
$\Omega$ と校正 & 将来の実現リターンで不確かさを決める & 禁止。X03 の校正は、保有期間が終わった過去 36 か月だけ \\
\end{longtable}

\section{LLM の先読み}

LLM の先読みは、ほかのデータの先読みより見つけにくい。gpt-4o-mini に 2022 年の決算発表を読ませると、モデルはその後の株価の動き、ニュース、アナリストの評価を学習のデータで読んでいる可能性がある。点数が反応をよく説明しても、それが「文章を読めた」からなのか「結果を覚えていた」からなのかを区別できない。

\begin{itemize}
\item \textbf{締め切りの後だけを使う}：gpt-4o-mini の点数は、2023-11 以降の発表だけで評価した（標本 A）。ただし、締め切りの直後は、モデルの更新でその後の情報が混ざる可能性が残る（\code{gpt-4o-mini-2024-07-18} は 2024 年 7 月の版）。
\item \textbf{版の決まったモデル}：ChronoGPT の版 $y$ は、$y$ 年 12 月 31 日までの文章だけで学習されている。判断日の前の年末の版を使えば、先読みはありえない。実験 5（2016〜2023 年）と実験 11（版 2015 で 2016〜2026 年）は、gpt-4o-mini の結果が記憶でないことの傍証になった。
\item \textbf{不利な比較}：実験 11 では、ChronoGPT を版 2015 に固定し、2026 年の発表まで読ませた。10 年前の言葉しか知らないモデルが、毎四半期学び直す TF-IDF と、2019 年までの文章を知る FinBERT を超えたので、LLM に不利な条件での勝ちである。
\end{itemize}

\section{生存者バイアスの大きさ}

yfinance には、上場を廃止した銘柄の株価がないことが多い。

\begin{itemize}
\item 総合報告書 1 の期間（2024〜2026 年）では、構成銘柄のうち 23 銘柄の株価がなく、28 銘柄を除いた。大半は買収された銘柄で、破綻した銀行（SVB Financial、First Republic）を含む。
\item DB を 2015 年まで広げると、期間中に指数に入った銘柄のうち 157 銘柄の yfinance の株価がなかった。別のデータの提供元（EODHD）で確かめると、そのうち 134 銘柄は上場廃止の一覧にあった。EODHD で取れる期間は約 12 年分までで、30 年前までさかのぼる取得は断念した。
\end{itemize}

この研究は、各判断日の構成銘柄を使い、株価がない銘柄だけを除いた。総合報告書 1 では、再現した各日の銘柄数が 503〜505 で実際の \SP{} と合うことを確かめたうえで、除いた銘柄による生存者バイアスの影響を「大」と評価した。戦略どうしの比較（LLM の BLM と時価総額加重、本物の点数と混ぜた点数）では、同じ銘柄の集合を使うので、比較そのものは公平である。ただし、成績の水準（時価総額加重そのもののリターン）は、実際の指数より良く出ている可能性がある。

\begin{projbox}
\begin{itemize}
\item \textbf{DB の時刻}：決算発表は受理時刻（UTC）、財務諸表は提出日、レポートは PDF の作成日、点数は付けた日時と使ったモデルの版。
\item \textbf{テスト}：時刻の変換（夏時間）、時間帯の境界、週末と祝日の day 0、受理日より後に提出された値を標準偏差に使わないことを、自動のテストで確かめた（第 VI 部）。
\item \textbf{実験の約束}：将来の実現リターンを参照して $\Omega$ を計算するなどの漏れは禁止。
\end{itemize}
\end{projbox}
